%! Designing Experiments — Unit 3, Lesson 3.4 — Homework (Answer Key)
% MIRROR of ../homework/main.tex: -key for -boxes, the header, and the correct MC option wrapped in
% \ans{}. Every \answerspace and \gradedblank line is identical; key mode prints each display.
% NO teachernote here — teacher prose lives in the lesson plan.
% GRADUAL RELEASE: this is the lesson's SCORED practice and the LAST component in the packet.
% Paper homework is the DEFAULT for this course; DeltaMath is an occasional override that
% replaces this page. Started in the last eight minutes
% of the period and finished after two study halls. The AP Practice component before it is extra credit,
% scored 1 to 5.
% THIRD CONTEXT: the notes teach on gum and memory, a sports drink, and a shoe sole; Guided
% Practice runs on a zinc lozenge; these items are on corn plots, pillows, and older dogs.
% Part A (corn on two kinds of field): A1 units/treatments/response, A2 why block, A3 THE CRUX
% restated — "already blocked, so plant Early Gold nearest the road". Part B: B1 the AP-style
% multiple choice (matched pairs, coin flip for order) with a required justification, B2 a
% placebo and blinding for a dog joint chew. 10 points.
% \answerspace{H}{} reserves exactly H; each \gradedblank inside it is BYTE-IDENTICAL in the key,
% where key mode prints its display. Keep every \boxguard identical in the blank and the key.
\documentclass[12pt]{article}
\usepackage{apstats-article}
\usepackage{apstats-key}   % pulls in -boxes; do NOT also load -boxes
\usepackage{grading}
\gradingassignment{apstats-u03-l04-hw}
\gradingcourse{advanced_statistics}
\newcommand{\answerspace}[2]{\par\nopagebreak\noindent\begin{minipage}[t][#1][t]{\linewidth}%
  \color{keyred}\bfseries #2\end{minipage}\par}

\begin{document}

\pageheader{Unit 3, Lesson 3.4}{Homework --- Answer Key}
{\small\itshape Write each final answer in its box. Only what is inside a box is graded.}\par
% NAMESTRIP: no \namedateperiod here — the name row lives on the cover only.

\vspace{0.04in}

% Authored BYTE-IDENTICALLY in homework/ and homework_key/.
\boxguard[8]
\begin{remindbox}
\small
\textbf{This is your graded homework.} It is scored 1 to 5 and due the first class after two study
halls. Start it now; ask about anything that stalls you before you leave.
\end{remindbox}

\vspace{0.04in}

\boxguard[14]
\begin{notesbox}{Part A --- Corn on Two Kinds of Field}
\interlinepenalty=10000 % keep each prompt whole when the box breaks
\small
An agronomist compares two corn varieties, \textbf{Early Gold} and \textbf{Prairie}, on 24
plots: 12 plots in the low, wet river bottom and 12 on the dry hillside. Soil moisture strongly
affects how much corn a plot produces. At harvest she measures the yield of each plot in
bushels per acre.
\begin{enumerate}[label=\textbf{A\arabic*.}, leftmargin=2.2em, itemsep=2pt]
  \item Name the experimental units, the treatments, and the response variable.
        \answerspace{2.1in}{\normalcolor\mdseries\hspace*{10.5pt}\gradedblank[id=hw1, type=rubric, points=2, display={Units: the 24 plots. Treatments: Early Gold and Prairie. Response: corn yield of each plot (bushels per acre).}]{\linewidth-10.5pt}{0.95cm}}
  \item She blocks on field: river bottom or hillside. Explain why this randomized block design
        is better here than a completely randomized design.
        \answerspace{2.1in}{\normalcolor\mdseries\hspace*{10.5pt}\gradedblank[id=hw2, type=rubric, points=2, display={Soil moisture affects yield. Comparing the varieties within each field removes the wet/dry difference; a completely randomized design could, by chance, put more of one variety in the wet field.}]{\linewidth-10.5pt}{0.95cm}}
  \item A farmhand says, ``The plots are already blocked by field, so we do not need to
        randomize. In each field, just plant Early Gold on the 6 plots nearest the road.''
        Respond.
        \answerspace{2.1in}{\normalcolor\mdseries\hspace*{10.5pt}\gradedblank[id=hw3, type=rubric, points=2, display={Disagree. Blocking does not replace random assignment. Plots near the road may differ (dust, drainage, sun), so Early Gold's yield would be mixed up with the road. Chance must pick the 6 plots in each field.}]{\linewidth-10.5pt}{0.95cm}}
\end{enumerate}
\end{notesbox}

\vspace{0.04in}

\boxguard[30]
\begin{notesbox}{Part B --- Pairs and Placebos}
\interlinepenalty=10000 % keep each prompt whole when the box breaks
\small
\begin{enumerate}[label=\textbf{B\arabic*.}, leftmargin=2.2em, itemsep=2pt]
  \item \textbf{AP-style.} A company wants to know whether its new pillow raises a
        sleep-quality score compared with a standard pillow. It has 40 volunteers. Which plan
        is a matched pairs design that uses chance correctly? Write its letter in the Letter
        box, then justify it in a sentence.
        \begin{enumerate}[label=\textbf{(\Alph*)}, leftmargin=2.2em, itemsep=0pt]
          \item The company picks 20 volunteers for the new pillow; the rest use the standard.
          \item Each volunteer sleeps a week on each pillow, always the standard one first.
          \item \ans{Each volunteer sleeps a week on each pillow; a coin flip decides which is first.}
          \item Each volunteer picks a pillow, and the two groups' scores are compared.
          \item 20 volunteers are randomly assigned to each pillow for one week.
                \hfill\gradedright{\textit{Letter:}~\gradedblank[id=hw4, type=choice, points=1, answer=C]{0.45in}{0.24in}}\rule[-0.2in]{0pt}{0.48in}
        \end{enumerate}
        \answerspace{2.1in}{\normalcolor\mdseries\hspace*{10.5pt}\gradedblank[id=hw4j, type=rubric, points=1, display={(C): each volunteer is compared with themselves on both pillows, and chance decides the order. (B) pairs but never randomizes the order.}]{\linewidth-10.5pt}{0.7cm}}
  \item A vet tests whether a joint chew helps older dogs move more easily. Owners rate their
        dog's mobility after 8 weeks. Explain why the control group should get a look-alike
        chew with no supplement, and say who should be kept blind.
        \answerspace{2.1in}{\normalcolor\mdseries\hspace*{10.5pt}\gradedblank[id=hw5, type=rubric, points=2, display={An owner who knows the dog got the supplement may expect, and report, better mobility; a placebo chew separates that from the supplement's effect. The owners and the vet staff who record ratings should not know which chew each dog got (double-blind).}]{\linewidth-10.5pt}{0.95cm}}
\end{enumerate}
\end{notesbox}

\vspace{0.04in}

% The next-lesson preview closes the packet, so it lives here rather than in AP Practice.
% A one-line, unbreakable box: guard it by its own height so it moves whole.
\boxguard[3]
\begin{spiralbox}
\small
\textbf{Next: Inference and Experiments} --- when a difference between treatment groups is too
big to be chance alone, and what an experiment lets us conclude.
\end{spiralbox}

\end{document}
